%% ============================================================
%%  Sistema de Historial Médico Familiar Centralizado
%%  con Alertas de Medicamentos
%%  Proyecto Integrador II — Junio 2026
%% ============================================================
\documentclass[12pt, letterpaper]{report}

%% ---- Codificación y lenguaje ----
\usepackage[utf8]{inputenc}
\usepackage[T1]{fontenc}
\usepackage[spanish, activeacute]{babel}
\addto\captionsspanish{%
  \renewcommand{\contentsname}{Tabla de Contenidos}%
  \renewcommand{\listfigurename}{Lista de Figuras}%
  \renewcommand{\listtablename}{Lista de Tablas}%
  \renewcommand{\bibname}{Referencias}%
  \renewcommand{\chaptername}{Capítulo}%
}

%% ---- Márgenes ----
\usepackage[top=2.54cm, bottom=2.54cm, left=3cm, right=2.54cm, headheight=15pt]{geometry}

%% ---- Interlineado ----
\usepackage{setspace}
\onehalfspacing

%% ---- Fuente ----
\usepackage{times}

%% ---- Hipervínculos ----
\usepackage[hidelinks]{hyperref}
\usepackage{url}

%% ---- Bibliografía ----
\usepackage[numbers, sort&compress]{natbib}

%% ---- Encabezados ----
\usepackage{fancyhdr}
\pagestyle{fancy}
\fancyhf{}
\fancyhead[R]{\small\nouppercase{\leftmark}}
\fancyfoot[C]{\thepage}
\renewcommand{\headrulewidth}{0.4pt}

%% ---- Formato de capítulos ----
\usepackage{titlesec}
\titleformat{\chapter}[display]
  {\normalfont\Large\bfseries\centering}
  {CAPÍTULO \thechapter}{10pt}{\large}
\titlespacing*{\chapter}{0pt}{30pt}{20pt}

%% ---- Listas ----
\usepackage{enumitem}

%% ============================================================
\begin{document}

%% ---- Portada ----
\begin{titlepage}
  \centering
  \vspace*{1cm}
  {\Large\bfseries SISTEMA DE HISTORIAL MÉDICO FAMILIAR\\[6pt]
  CENTRALIZADO CON ALERTAS DE MEDICAMENTOS\par}
  \vspace{1.5cm}
  {\large Por\par}
  \vspace{0.6cm}
  {\large Rodriguez Alvarado Oliver Rogelio\par}
  {\large Aguilar Hernández Jesús Alejandro \par}
  {\large De León Ruiz Hugo Guillermo\par}
  \vspace{1.5cm}
  {\large Presentado como parte de\par}
  {\large los requisitos necesarios para\par}
  {\large completar el Proyecto en Equipo 3 de\par}
  \vspace{1.5cm}
  {\large PROYECTO INTEGRADOR II\par}
  \vspace{1cm}
  {\large Junio 2026\par}
  \vspace{1cm}
  {\large Prof.\ Dr.\ Marco Aurelio Nuño Maganda\par}
  \vfill
\end{titlepage}

%% ---- Páginas preliminares ----
\pagenumbering{roman}
\setcounter{page}{2}

%% ============================================================
%%  ABSTRACT - supusimos que en español tambien
%% ============================================================
\chapter*{Abstract}
\addcontentsline{toc}{chapter}{Abstract}

Los sistemas de gestión de registros médicos electrónicos han demostrado ser herramientas
fundamentales para mejorar la calidad de la atención en salud, especialmente donde la
continuidad del historial clínico es crítica para la toma de decisiones~\cite{tsai2020effects}.
Sin embargo, en contextos familiares y de atención primaria persiste la necesidad de contar
con soluciones accesibles, centralizadas e interoperables que integren alertas activas sobre
el uso de medicamentos~\cite{gamache2018personal}.

El presente documento describe el diseño y desarrollo de un \textbf{Sistema de Historial
Médico Familiar Centralizado con Alertas de Medicamentos}, implementado sobre una
arquitectura web basada en PHP y MySQL, desplegada en la plataforma InfinityFree.
El sistema contempla módulos de autenticación con roles de usuario, expediente del
paciente, alta médica y notificaciones farmacológicas.

La metodología adoptada comprende cinco fases: planificación e investigación, diseño y
arquitectura, desarrollo del sistema, pruebas y validación, y documentación, con hitos de
entrega formales en cada fase.


%% ============================================================
%%  TABLA DE CONTENIDOS
%% ============================================================
\tableofcontents
\newpage

\listoftables
\addcontentsline{toc}{chapter}{Lista de Tablas} %% POR EL MOMENTO VACIA HASTA QUE LLEGUEMOS AL CAP 4
\newpage

\listoffigures
\addcontentsline{toc}{chapter}{Lista de Figuras} %% TAMBIEN LO MISMO DE ARRIBA
\newpage

%% ---- Cuerpo del documento ----
\pagenumbering{arabic}
\setcounter{page}{1}

%% ============================================================
%%  CAPÍTULO 1: INTRODUCCION
%% ============================================================
\chapter{INTRODUCCIÓN}

\section{Antecedentes}

\subsection{El Registro Médico Electrónico como Pilar de la Salud Digital}

La gestión de la información clínica ha experimentado una transformación profunda desde
la adopción de los registros médicos electrónicos (EHR). Estos sistemas permiten
almacenar, organizar y recuperar datos de salud de forma estructurada, reemplazando
los expedientes en papel que predominaron durante décadas en la práctica
médica~\cite{tsai2020effects}. Un EHR bien implementado incluye información
demográfica, diagnósticos, prescripciones, resultados de laboratorio e historiales de
visitas, conformando un repositorio integral del estado de salud del paciente~\cite{hayrinen2008definition}.

En el contexto familiar, la necesidad de contar con un historial centralizado se vuelve
particularmente relevante. Tang et al.~\cite{tang2006personal} señalan que los registros
personales de salud (PHR) ofrecen ventajas al empoderar al paciente y a su familia en
la gestión proactiva de su propio bienestar, reduciendo errores derivados de información
incompleta o desactualizada.

\subsection{El Problema de la Adherencia Farmacológica y las Alertas de Medicamentos}

Uno de los desafíos más persistentes en la atención sanitaria es la falta de adherencia
al tratamiento farmacológico. Este fenómeno provoca complicaciones evitables,
rehospitalizaciones y un incremento en los costos del sistema de
salud~\cite{gokoel2020medication}. Almomani et al.~\cite{almomani2023adherence}
encontraron que incluso en el contexto de antibióticos de corto plazo, una proporción
significativa de pacientes no completa el esquema prescrito.

Las intervenciones digitales orientadas a mejorar la adherencia han mostrado resultados
prometedores. Palmer et al.~\cite{palmer2021mobile} documentaron que las intervenciones
basadas en teléfono móvil reducen el incumplimiento terapéutico en adultos con riesgo
cardiovascular. Asimismo, Cao et al.~\cite{cao2022mhealth} demostraron que las
aplicaciones de mHealth con recordatorios y retroalimentación personalizada mejoran el
control de enfermedades crónicas como la hipertensión.

\subsection{Integración Tecnológica en Salud: Tendencias y Brechas}

La convergencia entre la informática médica, el desarrollo web y la salud móvil ha
impulsado una nueva generación de herramientas clínicas. La Organización Mundial de
la Salud~\cite{world2022consolidated} ha publicado guías para la implementación de
telemedicina que enfatizan la importancia de sistemas accesibles, seguros e
interoperables. Ghadi et al.~\cite{ghadi2025integration} destacan el creciente papel
de la inteligencia artificial y los dispositivos portátiles en el monitoreo remoto de
pacientes. Sin embargo, persiste una brecha importante entre las soluciones comerciales
de alto costo y la disponibilidad real de herramientas para entornos con recursos
limitados~\cite{roehrs2017omniphr}.

\section{Objetivos}

\subsection{Objetivo General}

Diseñar y desarrollar un sistema web de historial médico familiar centralizado con
alertas de medicamentos, que permita registrar, consultar y gestionar expedientes
clínicos de múltiples integrantes de un grupo familiar, integrando notificaciones
sobre el uso correcto y oportuno de los medicamentos prescritos.

\subsection{Objetivos Específicos}

\begin{enumerate}[label=\arabic*.]
  \item Investigar y analizar al menos 20 sistemas EHR/PHR existentes y artículos
        científicos relevantes para fundamentar el diseño del sistema propuesto.
  \item Levantar los requerimientos funcionales del sistema e identificar los roles
        de usuario necesarios para su operación segura.
  \item Modelar la base de datos relacional en MySQL que soporte el almacenamiento
        de expedientes, prescripciones y eventos de alerta.
  \item Diseñar las interfaces de usuario mediante wireframes y mockups que guíen
        la implementación del frontend.
  \item Implementar los módulos de autenticación, expediente del paciente, alta
        médica y alertas de medicamentos sobre la arquitectura PHP~+~MySQL.
  \item Desplegar el sistema en la plataforma InfinityFree y validarlo mediante
        pruebas unitarias, de integración y de usabilidad con usuarios piloto.
  \item Documentar el sistema mediante un manual de usuario y documentación técnica
        que faciliten su mantenimiento y futura extensión.
\end{enumerate}

\section{Relevancia e Importancia}

La relevancia de este proyecto se articula en tres dimensiones: clínica, tecnológica
y social. Desde la perspectiva clínica, un sistema centralizado de historial familiar
reduce el riesgo de errores médicos derivados de información fragmentada o
desactualizada~\cite{hayrinen2008definition}. Las alertas de medicamentos contribuyen
a prevenir interacciones farmacológicas peligrosas y a mejorar la adherencia
terapéutica en padecimientos crónicos~\cite{etminani2020behavior, arshed2024effectiveness}.

En el plano tecnológico, el proyecto demuestra la viabilidad de implementar
funcionalidades avanzadas de gestión clínica sobre un stack accesible
(PHP~+~MySQL~+~HTML/CSS/JS). Alsyouf et al.~\cite{alsyouf2023use} señalan que la
usabilidad percibida y la confianza en la privacidad de los datos son determinantes
clave en la adopción de sistemas PHR; ambos aspectos son abordados en el diseño del
presente sistema.

Finalmente, en términos sociales, el proyecto contribuye a democratizar el acceso a
herramientas de salud digital para familias que no cuentan con acceso a soluciones
EHR comerciales. La OMS~\cite{world2022consolidated} identifica este tipo de
soluciones como palancas estratégicas para la cobertura sanitaria universal.

\section{Estructura del Informe}

El presente informe se organiza siguiendo una progresión lógica desde el contexto
general del problema hasta los fundamentos técnicos que sustentan el proyecto:

\medskip
\noindent\textbf{Revisión de Literatura.} Se presenta un análisis crítico de los
sistemas EHR/PHR existentes, las intervenciones digitales para adherencia
farmacológica y las tendencias en salud digital, con base en las referencias
bibliográficas pertinentes. Se identifican las brechas que motivan el desarrollo
del presente sistema.

\medskip
\noindent\textbf{Enfoque Técnico.} Se describe la metodología de desarrollo adoptada,
incluyendo el levantamiento de requerimientos, el modelado de la base de datos, el
diseño de interfaces, la arquitectura del sistema y las estrategias de prueba,
estructuradas en las cinco fases del plan de proyecto.

\medskip
\noindent\textbf{Resultados.} Se reportan los resultados de cada fase de desarrollo,
desde la configuración del entorno hasta las pruebas de validación con usuarios
piloto, apoyándose en métricas de desempeño y matrices de prueba.

\medskip
\noindent\textbf{Discusión.} Se interpretan los resultados en relación con la
literatura existente, se destacan las contribuciones del proyecto y se analizan
sus limitaciones y las áreas de mejora identificadas.

\medskip
\noindent\textbf{Conclusión.} Se sintetizan los hallazgos principales, se evalúa
el grado de cumplimiento de los objetivos y se proponen líneas de trabajo futuro.

%% ============================================================
%%  CAPÍTULO 2: REVISIÓN DE LITERATURA
%% ============================================================
\chapter{REVISIÓN DE LITERATURA}

\section{Registros Médicos Electrónicos y Registros Personales de Salud}

Los registros médicos electrónicos (EHR) y los registros personales de salud (PHR)
constituyen los fundamentos conceptuales del presente proyecto. Häyrinenetal.~\cite{hayrinen2008definition}
ofrecen una definición ampliamente citada: un EHR es un repositorio longitudinal de
información de salud de un paciente, generado en uno o más encuentros en cualquier
entorno de atención, cuyo objetivo principal es el soporte a la atención médica continua
y la toma de decisiones. En contraste, los PHR otorgan al propio individuo el control
sobre sus datos, promoviendo un modelo de atención centrado en el
paciente~\cite{tang2006personal}.

Roehrs et al.~\cite{roehrs2017omniphr} desarrollaron OmniPHR, una arquitectura
distribuida que integra registros personales de salud mediante nodos interoperables,
demostrando que es posible construir sistemas robustos sin infraestructura centralizada
costosa. Por su parte, Tsai et al.~\cite{tsai2020effects} realizaron una revisión de
alcance sobre los efectos de la implementación de EHR e identificaron barreras
recurrentes: resistencia del personal, costos elevados, problemas de usabilidad y
preocupaciones sobre privacidad.

\section{Adherencia Farmacológica e Intervenciones Digitales}

La adherencia al tratamiento farmacológico es un problema global de salud pública.
Gokoel et al.~\cite{gokoel2020medication} documentaron tasas de no adherencia
superiores al 30\% en pacientes postrasplante, con consecuencias directas sobre la
supervivencia del injerto. En poblaciones con enfermedades crónicas como hipertensión,
el fenómeno es igualmente prevalente~\cite{etminani2020behavior}.

Las intervenciones basadas en tecnologías digitales han demostrado ser efectivas para
mitigar esta problemática. Palmer et al.~\cite{palmer2021mobile} concluyeron, en una
revisión Cochrane, que las intervenciones por teléfono móvil mejoran la adherencia a
medicamentos cardioprotectores. Cao et al.~\cite{cao2022mhealth} identificaron que los
componentes de mayor impacto son la personalización de los mensajes y la monitorización
continua. Arshed et al.~\cite{arshed2024effectiveness} evaluaron un paquete de
intervención mHealth multifactorial en un país de ingreso medio-bajo, encontrando
mejoras sustanciales tanto en adherencia como en control de la presión arterial.

Chandler et al.~\cite{chandler2019impact} demostraron que la adaptación cultural de
los programas digitales de autogestión de medicamentos es un factor clave para su
efectividad en poblaciones hispanas, aspecto especialmente relevante en el contexto
latinoamericano en el que se inscribe este proyecto.

\section{Aplicaciones Móviles y Sistemas de Alerta Farmacológica}

Okati et al.~\cite{okati2025mobile} revisaron aplicaciones disponibles en tiendas de
aplicaciones para la desprescripción, hallando que la mayoría carece de bases de
evidencia sólidas y que la interoperabilidad con los registros clínicos institucionales
es escasa. Eaton et al.~\cite{eaton2024user} analizaron el compromiso del usuario con
intervenciones mHealth en personas con enfermedades crónicas, concluyendo que la
usabilidad y la personalización son los predictores más robustos del uso sostenido.

Zavaleta-Monestel et al.~\cite{zavaleta2025artificial} presentaron una revisión
actualizada sobre herramientas de inteligencia artificial para mejorar la adherencia
en enfermedades crónicas no transmisibles, identificando el uso de chatbots y
algoritmos de predicción de riesgo como las aplicaciones más prometedoras. Birkun y
Gautam~\cite{birkun2024large} exploraron el uso de modelos de lenguaje de gran escala
como fuentes de consejo médico, señalando tanto el potencial como los riesgos de estas
tecnologías sin supervisión clínica.

\section{Salud Digital, Telemedicina y Seguridad de Datos}

Pool et al.~\cite{pool2024systematic} realizaron un análisis sistemático de fallos en
la protección de datos de salud personales, identificando como vulnerabilidades
frecuentes la autenticación débil, el cifrado insuficiente y la falta de auditorías
internas. Alsyouf et al.~\cite{alsyouf2023use} aplicaron el modelo de aceptación
tecnológica (TAM) para predecir el uso de sistemas PHR, encontrando que la percepción
de seguridad y privacidad es el factor de mayor peso en la intención de uso.

La OMS~\cite{world2022consolidated} enfatiza que cualquier sistema de telemedicina o
salud digital debe cumplir estándares de confidencialidad, integridad y disponibilidad
de la información. Ringeval et al.~\cite{ringeval2025advancing} analizaron las
aplicaciones de los gemelos digitales en salud, señalando que la integración de datos
clínicos longitudinales es el punto de partida para modelos predictivos avanzados.

\section{Sistemas de Gestión de Salud en Contextos de Recursos Limitados}

Katz et al.~\cite{katz2024digital} revisaron las intervenciones de salud digital para
el manejo de hipertensión en poblaciones con disparidades de salud, concluyendo que
las soluciones digitales pueden reducir brechas de acceso cuando se diseñan con
criterios de equidad. Stevenson et al.~\cite{stevenson2019ehealth} evaluaron
intervenciones eHealth en enfermedad renal crónica, encontrando beneficios en el
seguimiento de la adherencia.

Tong et al.~\cite{tong2022digital} y Duettmann et al.~\cite{duettmann2021ehealth}
abordaron la dimensión de la alianza terapéutica digital y la adopción de herramientas
eHealth en trasplante, respectivamente, destacando que la confianza del usuario en el
sistema es determinante del éxito de implementación. Gonzales et al.~\cite{gonzales2021pharmacist}
demostraron que las intervenciones lideradas por farmacéuticos a través de plataformas
móviles mejoran la seguridad en el uso de medicamentos, evidenciando la importancia
de integrar el componente farmacológico en los sistemas de gestión clínica. En una
línea similar, Abasi et al.~\cite{abasi2021effectiveness} y Duncan et al.~\cite{duncan2018systematic}
documentaron la efectividad de aplicaciones de autogestión y de intervenciones para
adherencia a inmunosupresores, respectivamente, reforzando la evidencia a favor de
soluciones digitales integradas. Por su parte, Huuskes et al.~\cite{huuskes2021kidney}
exploraron las perspectivas de pacientes trasplantados sobre la telehealth durante la
pandemia de COVID-19, destacando la aceptación creciente de este tipo de
herramientas. Finalmente, Campbell et al.~\cite{campbell2022interventions} analizaron
intervenciones para mejorar la alfabetización en salud en pacientes con enfermedad
renal crónica, subrayando que la comprensión del propio historial médico es un
habilitador clave de la adherencia terapéutica.

\section{Brecha en la Investigación Existente}

A pesar del volumen considerable de literatura sobre EHR, PHR e intervenciones
digitales para la adherencia farmacológica, se identifican brechas relevantes que
justifican el presente proyecto. En primer lugar, la mayoría de los sistemas revisados
han sido desarrollados para contextos de alta disponibilidad tecnológica, con
infraestructuras no replicables en entornos con recursos
limitados~\cite{roehrs2017omniphr, tsai2020effects}. En segundo lugar, la integración
de alertas farmacológicas dentro de un sistema de historial familiar centralizado es
un área con escasa producción científica~\cite{okati2025mobile}. Finalmente, la
evidencia sobre la efectividad de soluciones construidas con tecnologías web estándar
(PHP, MySQL) desplegadas en plataformas de alojamiento gratuito es prácticamente
inexistente, lo que otorga al presente proyecto un carácter exploratorio y de
contribución metodológica.

%% ============================================================
%%  REFERENCIAS
%% ============================================================
\newpage
\bibliographystyle{plainnat}
\bibliography{references}
\addcontentsline{toc}{chapter}{Referencias}

\end{document}
